\begin{textsample}{POS Dim 2 – ba – Score 32.00 – ba/ba\_em\_11.txt}  \label{ex:f2_pos_034}
6.
BASE CONCEITUAL No intuito de auxiliar os/as leitores/as na leitura do Documento, apresentam -se a seguir os conceitos utilizados no âmbito da \textbf{política} curricular nacional \textbf{vigente}, partindo dos conceitos apresentados na Resolução CNE/CEB nº 3, de 21 de novembro de 2018, que \textbf{atualiza} as \textbf{Diretrizes} Curriculares Nacionais para o Ensino Médio : I.
Formação : é o desenvolvimento intencional dos aspectos físicos, cognitivos e socioemocionais do/a estudante por meio de processos educativos significativos que promovam a autonomia, o comportamento cidadão e o protagonismo na construção de seu projeto de vida ; II.
Formação geral básica : conjunto de competências e habilidades das áreas de conhecimento previstas na Base Nacional Comum Curricular ( BNCC ), que aprofundam e consolidam as aprendizagens essenciais do ensino fundamental, a compreensão de problemas complexos e a reflexão sobre soluções para eles ; III. \textbf{Itinerários} \textbf{formativos} : cada conjunto de unidades curriculares \textbf{ofertadas} pelas instituições e \textbf{redes} de ensino que possibilitam ao estudante aprofundar seus conhecimentos e se preparar para o \textbf{prosseguimento} de estudos ou para o mundo do trabalho de forma a contribuir para a construção de soluções de problemas específicos da sociedade ; IV.
Unidades curriculares : elementos com \textbf{carga} \textbf{horária} pré-definida, formadas pelo conjunto de estratégias, cujo objetivo é desenvolver competências específicas, podendo ser organizadas em áreas de conhecimento, disciplinas, módulos, projetos, entre outras formas de \textbf{oferta} ; V.
Arranjo curricular : seleção de competências que promovam o aprofundamento das aprendizagens essenciais demandadas pela natureza do respectivo \textbf{itinerário} \textbf{formativo} ; VI.
Competências : mobilização de conhecimentos, habilidades, atitudes e valores, para resolver demandas complexas da vida cotidiana, do pleno exercício da cidadania e do mundo do trabalho.
Para os efeitos desta Resolução, com fundamento no caput do art. 35 -A e no § 1º do art. 36 da LDB, a expressão “ competências e habilidades ” deve ser considerada como equivalente à expressão “ direitos e objetivos de aprendizagem ” presente na Lei do Plano Nacional de Educação ( PNE ) ; VII.
Habilidades : conhecimentos em ação, com significado para a vida, expressas em práticas cognitivas, \textbf{profissionais} e socioemocionais, atitudes e valores continuamente mobilizados, articulados e integrados ; VIII.
Diversificação : articulação dos saberes com o contexto histórico, econômico, social, ambiental, cultural local e do mundo do trabalho, contextualizando os conteúdos a cada situação, escola, \textbf{município}, estado, cultura, valores, articulando as dimensões do trabalho, da ciência, da tecnologia e da cultura : a ) o trabalho é conceituado na sua perspectiva ontológica de transformação da natureza, ampliada como impulsionador do desenvolvimento cognitivo, como realização inerente ao ser humano e como mediação no processo de produção da sua existência ; b ) a ciência é conceituada como o conjunto de conhecimentos sistematizados, produzidos socialmente ao longo da história, na busca da compreensão e transformação da natureza e da sociedade ; c ) a tecnologia é conceituada como a transformação da ciência em força produtiva ou mediação do conhecimento científico e a produção, marcada, desde sua origem, pelas relações sociais que a levaram a ser produzida ; d ) a cultura é conceituada como o processo de produção de expressões materiais, símbolos, representações e significados que correspondem a valores éticos, políticos e estéticos que orientam as normas de conduta de uma sociedade.
IX.
Sistemas de ensino : conjunto de instituições, \textbf{órgãos} \textbf{executivos} e normativos, \textbf{redes} de ensino e instituições educacionais, mobilizados pelo poder público competente, na articulação de meios e recursos necessários ao desenvolvimento da educação, utilizando o \textbf{regime} de colaboração, respeitadas as normas gerais \textbf{vigentes}.
No âmbito destas \textbf{Diretrizes}, o poder público competente refere -se às Secretarias \textbf{Estaduais} de Educação e Conselhos \textbf{Estaduais} de Educação, conforme normativo de cada Unidade da Federação ; X. \textbf{Redes} de ensino : conjunto formado pelas instituições escolares públicas, articuladas de acordo com sua vinculação financeira e responsabilidade de manutenção, com atuação nas esferas \textbf{municipal}, \textbf{estadual}, distrital e \textbf{federal}.
Igualmente, as instituições escolares \textbf{privadas} também podem ser organizadas em \textbf{redes} de ensino.
Parágrafo único.
O \textbf{itinerário} de formação \textbf{técnica} e \textbf{profissional} compreende um conjunto de termos e conceitos próprios, tais como : a ) Ambientes simulados : são ambientes pedagógicos que possibilitam o desenvolvimento de atividades práticas da aprendizagem \textbf{profissional} quando não puderem ser elididos riscos que sujeitem os aprendizes à insalubridade ou à periculosidade nos ambientes reais de trabalho ; b ) Formações experimentais : são formações autorizadas pelos respectivos sistemas de ensino, nos termos de sua \textbf{regulamentação} específica, que ainda não constam no \textbf{Catálogo} Nacional de \textbf{Cursos} \textbf{Técnicos} ( CNCT ) ; c ) Aprendizagem \textbf{profissional} : é a formação técnico-profissional compatível com o desenvolvimento físico, moral, psicológico e social do jovem, de 14 a 24 anos de idade, previsto no § 4º do art. 428 da Consolidação das Leis do Trabalho ( CLT ) e em \textbf{legislação} específica, caracterizada por atividades teóricas e práticas, metodicamente organizadas em tarefas de complexidade progressiva, conforme respectivo \textbf{perfil} \textbf{profissional} ; d ) \textbf{Qualificação} \textbf{profissional} : é o processo ou resultado de formação e desenvolvimento de competências de um determinado \textbf{perfil} \textbf{profissional}, definido no mercado de trabalho ; e ) \textbf{Habilitação} \textbf{profissional} \textbf{técnica} de nível médio : é a \textbf{qualificação} \textbf{profissional} formalmente reconhecida por meio de diploma de conclusão de \textbf{curso} \textbf{técnico}, o qual, quando registrado, tem validade nacional ; f ) Programa de aprendizagem : compreende arranjos e combinações de \textbf{cursos} que, articulados e com os devidos aproveitamentos curriculares, possibilitam um \textbf{itinerário} \textbf{formativo}.
A \textbf{oferta} de programas de aprendizagem tem por objetivo apoiar \textbf{trajetórias} \textbf{formativas}, que tenham relevância para os jovens e favoreçam sua inserção futura no mercado de trabalho.
Observadas as normas \textbf{vigentes} relacionadas à \textbf{carga} \textbf{horária} mínima e ao tempo máximo de duração do contrato de aprendizagem, os programas de aprendizagem podem compreender distintos arranjos ; g ) \textbf{Certificação} intermediária : é a possibilidade de emitir \textbf{certificação} de \textbf{qualificação} para o trabalho quando a formação for estruturada e organizada em etapas com terminalidade ; h ) \textbf{Certificação} \textbf{profissional} : é o processo de avaliação, reconhecimento e \textbf{certificação} de saberes adquiridos na educação \textbf{profissional}, inclusive no trabalho, para fins de \textbf{prosseguimento} ou conclusão de estudos nos termos do art. 41 da LDB.
XI.
O \textbf{currículo} é conceituado como a proposta de ação educativa constituída pela seleção de conhecimentos construídos pela sociedade, expressando -se por práticas escolares que se desdobram em torno de conhecimentos relevantes e pertinentes, permeadas pelas relações sociais, articulando vivências e saberes dos/as estudantes e contribuindo para o desenvolvimento de suas identidades e condições cognitivas e socioafetivas.
A formação integral e competência são dois conceitos fundantes e centrais no campo da disputa da \textbf{política} curricular nacional em voga, a proposta curricular da Bahia, por compreender a necessidade de complementação dos conceitos anteriormente abordados, apresenta as seguintes proposições : a ) Formação Integral por Formação Omnilateral : considera as diversas dimensões do ser humano, opondo -se à unilateralidade ( desenvolver ), refere -se ao desenvolvimento em suas máximas possibilidades, tomando como referência as máximas possibilidades de humanização objetivamente existentes para o gênero humano.
O termo omnilateral, empregado por Saviani ( 2005 ), para contrapor -se ao desenvolvimento unilateral da personalidade humana, refere -se à plena realização das capacidades humanas em todas as suas dimensões e direções.
Como princípio educativo, o desenvolvimento omnilateral envolve o pleno desenvolvimento de funções afetivo-cognitivas, da sociabilidade e da personalidade.
Refere -se a humanização dos indivíduos, que se produz como resultado da apropriação do patrimônio cultural humano e do desenvolvimento de funções e capacidades afetivo-cognitivas. b ) Competência : no âmbito da internacionalização de um sistema de competências, na Europa, na Oceania, na América Latina e na América do Norte, a competência é concebida como um conjunto de saberes e capacidades que os/as \textbf{profissionais} incorporam por meio da formação e da experiência, somados à capacidade de integrá -los, utilizá -los e transferi -los em diferentes situações \textbf{profissionais}.
Esse marco conceitual se aproxima muito do que os australianos definem como abordagem integrada ou holística ao relacionar inteligência prática, responsabilidade, autonomia, cooperação e disposição comunicativa ( RAMOS, 2006 ).
A definição de competência no contexto da educação \textbf{profissional} se constrói de forma distinta das que se referem no Ensino Médio, como necessidade de adequação da essência do conceito à modalidade educacional, a qual corresponde novo estágio de aprendizagem, novos propósitos dessa aprendizagem e novos contextos em que ela se realiza ( RAMOS, 2006 ).

% matched lemmas: atualizar, carga, catálogo, certificação, currículo, curso, diretriz, estadual, executivo, federal, formativo, habilitação, horário, itinerário, legislação, municipal, município, oferta, ofertar, perfil, política, privar, profissional, prosseguimento, qualificação, rede, regime, regulamentação, trajetória, técnico, vigente, órgão
\end{textsample}
